\clearpage
\VerdierPhysicalPageAnchor{0026}
多加性函子到复形上的自然延拓不一定与平移严格交换，但会“在同构意义下”
交换；在实际中，这些同构由一个符号确定。事实证明，在复合“平移函子”时，
不能相容地选取这些同构，因而会引入一个符号。第 1 节研究这个符号（另见
附录）。第 2 节定义多重复形、单复形、“关联单复形”函子与同伦等经典概念；
第 3 节开始研究复形态射的锥（“映射锥”或“映射柱”），这是后文的关键
构造。

\VerdierSemanticAnchor{intro-6-plan-derived-category}
阿贝尔范畴 $\mathcal A$ 的导出范畴 $\mathbf D(\mathcal A)$ 分两步构造。
第一步引入 $\mathcal A$ 的复形在同伦意义下所成的范畴
$\mathbf K(\mathcal A)$；然后在 $\mathbf K(\mathcal A)$ 中形式地逆转
拟同构，得到 $\mathbf D(\mathcal A)$。范畴 $\mathbf K(\mathcal A)$ 与
$\mathbf D(\mathcal A)$ 不一定是阿贝尔范畴，但它们带有额外结构，即由
锥构造定义的一族图式：\emph{特异三角}。对于这些范畴，特异三角扮演阿贝尔
范畴中正合列的角色。

\VerdierSemanticAnchor{intro-6-plan-chapter-two}
带有这一额外结构的加性范畴由 Puppe [2] 提出并研究，称为
\emph{三角范畴}；第二章从我们的观点研究它们。第 1 节陈述三角范畴的若干
初等结果。第 2 节在三角范畴的一般框架中证明：形式地逆转三角范畴
$\mathbf K(\mathcal A)$ 的拟同构这一问题，可由“分式演算”[16] 简洁解决。
第 3 节证明 Freyd [17] 的一个定理：三角范畴 $\mathcal D$ 可按普遍方式
嵌入阿贝尔范畴 $\mathrm A(\mathcal D)$。该范畴有一族同时为内射且为投射
的对象 $E$，使 $\mathrm A(\mathcal D)$ 的每个对象既同构于某个 $E$ 中
对象的子对象，又同构于某个 $E$ 中对象的商对象。由 $E$ 定义的
$\mathrm A(\mathcal D)$ 的全子范畴同构于 $\mathcal D$。最后，第 4 节
研究取值于三角范畴或阿贝尔范畴的“谱对象”；这里我们只是改写 [1] 中的
定义。

\VerdierSemanticAnchor{intro-6-plan-chapter-three}
第三章专门研究导出范畴本身。设 $X,Y$ 是阿贝尔范畴 $\mathcal A$ 的导出
范畴 $\mathbf D(\mathcal A)$ 中的两个对象。若 $\mathcal A$ 有足够多的
内射对象，且复形 $Y$ 的分量在足够小的次数上为零，则从 $X$ 到 $Y$ 的
次数为 $n$ 的态射可解释为
